\subsection{关系感知的可学习结构传播模块}

在多源证据驱动的结构传播式故障定位中，静态调用/依赖关系与动态覆盖关系为候选代码行提供了重要的结构约束。然而，真实工程场景中图结构往往并不可靠：一方面，覆盖与依赖边可能包含噪声或冗余连接；另一方面，固定图结构会把错误信号在局部邻域内放大，从而导致排序结果对边噪声与局部扰动过于敏感。为此，本节提出一种关系感知的可学习结构传播模块，将“边是否可信、是否需要被保留”作为可学习对象，而不是仅在固定图上做强制传播，以提升排序的稳健性与可解释性。

（1）候选结构构建与关系感知表示

首先以缺陷版本为基本样本单位构建候选异构图，保留覆盖、归属、依赖/调用等多关系边作为“候选结构集合”。其中，测试结点、代码行结点与方法结点构成不同类型的结点集合；覆盖关系与归属关系用于维持跨模态对齐通道，依赖/调用关系用于提供跨位置传播路径。在此基础上，为每一类关系引入关系类型嵌入，并对结点对\((i,j)\)构建关系感知表征\([h_i;h_j;e_r]\)，为后续结构学习提供输入。

（2）可学习边权与稀疏正则

对每条候选边\((i\!\rightarrow\!j,r)\)引入边权生成器，利用结点对表示预测边权：
\[
g_{ij}=\sigma\left(W_g [h_i;h_j;e_r]\right),
\]
其中\(\sigma(\cdot)\)为 Sigmoid。为避免结构被过度改写导致排序不稳定，本文采用“软门控”方式将学习到的边权与原始结构进行残差融合：
\[
\tilde{a}_{ij}=(1-\alpha)\cdot a_{ij}+\alpha\cdot g_{ij}\cdot a_{ij},
\]
其中\(\alpha\)控制结构学习强度。为了得到更稀疏、可解释的子图结构，本文对边权引入\(L_1\)稀疏正则，鼓励模型在候选边中自动筛除冗余或噪声连接。该机制在保留必要结构通道的同时，抑制了不可靠边对传播结果的干扰。

（3）结构感知传播与排序学习

在得到可学习的边权后，模型在异构图上执行关系感知的消息传递，将\(\tilde{a}_{ij}\)作为传播权重参与注意力归一化与邻域聚合；随后对代码行结点输出可疑度分数，并以缺陷版本为单位进行排序学习。通过“候选结构构建—边权学习—稀疏正则—结构感知传播”的流程，模型能够在保留覆盖与结构证据优势的同时，自动削弱噪声边的影响，从而在候选同质化场景下提升排序稳定性，并为后续分析提供更具可解释性的关键边分布。

该模块的设计为结构传播引入了“可学习结构选择”的能力，使模型不再依赖固定图结构的完备性与可靠性。在实验层面，可通过“稀疏度—性能曲线”“关系边敏感性”与“关键边类型分布”来验证结构学习的有效性与可解释性。接下来将结合整体框架，给出该模块的具体实现与实验设定。
